\chapter{Point-in-Time Data and the Biases}\label{ch:rs:point-in-time-data-and-the-biases}

On 3 October 2008 the US employment report said that the economy had lost 159\,000 jobs in September. Today's database says
451\,000. Between the two, the figure was revised six times, three of them by the annual benchmark revisions that the
statistical agency publishes each February. A macro signal backtested on
today's database would have traded on a number nobody had in October 2008, and would have seen a recession deeper and clearer
than the one anyone saw at the time. This chapter is about the difference between what is true now and what was known then:
\omterm{def:rs:point-in-time-data-and-the-biases:lookahead}{look-ahead bias} and the time stamps that prevent it, \omterm{def:rs:point-in-time-data-and-the-biases:survivorship}{survivorship bias}, revisions and \omterm{def:rs:point-in-time-data-and-the-biases:vintage}{restatements}, the storage that keeps
every vintage of a fact, and the \omterm{def:rs:point-in-time-data-and-the-biases:asof}{as-of join} that reads it. Its data are the real-time vintages of US payrolls and a
simulated stock market with delistings.

\section{Look-ahead bias and timestamp semantics}

\begin{definition}[Look-ahead bias, point-in-time data]\label{def:rs:point-in-time-data-and-the-biases:lookahead}
A backtest has \emph{look-ahead bias}\index{look-ahead bias} when a decision at time $t$ uses information that was not
available at $t$. \emph{Point-in-time data}\index{point-in-time data} are data stored so that, for any past instant, they
return exactly what was known at that instant: the values then published, the universe then listed, the identifiers then
used.
\end{definition}

Look-ahead rarely comes from a deliberate peek. It comes from a time stamp that means something other than what the code
assumes. A daily \omterm{def:rs:market-data-for-research:bar}{bar} stamped with its date is known only at the close, so a signal computed from it cannot trade at that
day's close. An earnings figure stamped with its fiscal quarter's end is known weeks later, on the filing date (chapter~11).
An index change stamped with its effective date was announced days before (One Quant Book 1, chapter 15), and the
announcement is when the trade starts. A trade stamped by the exchange arrives at the firm later, and the arrival time is the
one the firm could act on (Book 1, chapter 28). Each fact therefore needs two times.

\begin{definition}[Valid time, knowledge time, bitemporal data]\label{def:rs:point-in-time-data-and-the-biases:bitemporal}
The \emph{valid time}\index{valid time} of a value is the time it describes: September 2008 for September's payroll change, the
fourth quarter for a quarterly earnings figure. Its \emph{knowledge time}\index{knowledge time} is when it became available to
the user: the release, the filing, the vendor's delivery, the arrival of the message. Data are \emph{bitemporal}\index{bitemporal data}
when every value is stored with both, and a new value for the same valid time is added rather than overwriting the old one.
\end{definition}

\begin{omfigure}
\begin{tikzpicture}[font=\footnotesize,x=1.35cm,y=0.8cm]
\draw[->,>=Stealth] (0,0) -- (8.4,0) node[right] {valid time};
\draw[->,>=Stealth] (0,0) -- (0,5.4) node[above] {knowledge time};
\foreach \x/\lab in {1/Aug,2/Sep,3/Oct,4/Nov} \node[below] at (\x,0) {\lab};
\foreach \y/\lab in {1/Oct,2/Nov,3/Dec,4/{Feb, later}} \node[left] at (0,\y) {\lab};
\fill[omDef] (1,1) circle (2.2pt) (1,2) circle (2.2pt) (1,3) circle (2.2pt) (1,4) circle (2.2pt);
\fill[omDef] (2,1) circle (2.2pt) (2,2) circle (2.2pt) (2,3) circle (2.2pt) (2,4) circle (2.2pt);
\fill[omDef] (3,2) circle (2.2pt) (3,3) circle (2.2pt) (3,4) circle (2.2pt);
\fill[omDef] (4,3) circle (2.2pt) (4,4) circle (2.2pt);
\draw[omThm,very thick] (0.5,2) -- (4.6,2);
\node[omThm,anchor=west] at (4.8,2) {a snapshot: what was known in November};
\draw[black!60,very thick,dashed] (0.5,4) -- (4.6,4);
\node[black!70,anchor=west] at (4.8,4) {today's view: the latest vintage};
\draw[omExo,very thick] (2,0.55) -- (2,4.45);
\node[omExo,anchor=west] at (4.8,0.9) {one number's life: every vintage of September};
\draw[omExo,->,>=Stealth] (4.75,0.9) -- (2.12,1.4);
\end{tikzpicture}

\omcaption{\omterm{def:rs:point-in-time-data-and-the-biases:bitemporal}{Bitemporal data}. Each dot is a value stored for a month (\omterm{def:rs:point-in-time-data-and-the-biases:bitemporal}{valid time}) as published at a date (\omterm{def:rs:point-in-time-data-and-the-biases:bitemporal}{knowledge time}); a
month first appears the month after it and is revised later. A backtest must read horizontal cuts at its decision times,
never the top one; the vertical line through September is the life of one number, drawn in
\cref{fig:rs:point-in-time-data-and-the-biases:life}.}\label{fig:rs:point-in-time-data-and-the-biases:plane}
\end{omfigure}

\section{Survivorship bias}

\begin{definition}[Survivorship bias]\label{def:rs:point-in-time-data-and-the-biases:survivorship}
\emph{Survivorship bias}\index{survivorship bias} is the error of studying only the securities, funds or firms that exist at
the end of the sample: those that were delisted, liquidated or merged away are missing, and they are not missing at random.
\end{definition}

A stock universe taken from today's list of listed companies has already dropped every company that failed. A database of
funds built from the funds reporting today has dropped those that closed. The dropped names are, on average, the ones that did
badly, so their absence flatters every average. Brown, Goetzmann, Ibbotson and Ross (1992) showed how survivorship alone can
create apparent persistence in fund performance; for stocks, Shumway (1997) found that the delisting returns of stocks delisted
for poor performance were mostly missing from the standard US database, and that, where he could recover them from
over-the-counter prices, they averaged $-30\%$.

\begin{proposition}[How much the survivors overstate]\label{prop:rs:point-in-time-data-and-the-biases:survivors}
Pool all the name-months of a universe. Let a fraction $w$ of them belong to names that are delisted before the end of the
sample, with average return $r_F$, and the rest to names still listed at the end, with average return $r_S$. The pooled average
return of the whole universe is $(1 - w)r_S + wr_F$, so an average over the survivors overstates it by $w(r_S - r_F)$.
\end{proposition}

\begin{proof}
The pooled average is the weighted average of the two groups' averages, with weights equal to their shares of name-months.
\end{proof}

The identity says where the bias comes from: not only the delisting return, a single month, but every month a failing name
spends declining towards its delisting. The chapter's simulated market makes it concrete. Two thousand stocks follow random
walks in log price, with an 8\% drift and 40\% volatility a year; a stock whose cumulative log return falls below $-1.6$ (an 80\%
loss) is delisted with a return of $-30\%$ that month and replaced by a new listing. Over twenty years 2.2\% of names are delisted
each year. The equal-weighted universe, delistings included, returns 7.7\% a year; the universe of today's survivors,
backfilled over their own histories, returns 12.9\%: a bias of 5.2 points a year
(\cref{fig:rs:point-in-time-data-and-the-biases:surv}). In the proposition's terms, 20.9\% of the name-months belong to names
that later fail, with $r_F = -11.1\%$ a year against $r_S = 12.6\%$, so $w(r_S - r_F) = 5.0$ points; the remaining difference
comes from averaging month by month over a changing number of survivors rather than pooling. Setting the delisting return to
zero lowers the bias only to 4.5 points: the decline before delisting is most of it.

\begin{omfigure}
\begin{tikzpicture}
\begin{semilogyaxis}[width=10cm,height=5.2cm,xlabel={years},ylabel={growth of 1 (log scale)},
  tick label style={font=\small},label style={font=\small},xmin=0,xmax=20,ymin=0.8,ymax=16,grid=major,grid style={gray!25},
  ytick={1,2,4,8,16},yticklabels={1,2,4,8,16},legend columns=2,legend style={font=\footnotesize,at={(0.5,-0.32)},
  anchor=north,draw=none,fill=none,/tikz/every even column/.append style={column sep=8pt}},table/col sep=comma]
\addplot[omDef,very thick] table[x=year,y=full]{figdata/research/03-point-in-time-data-and-the-biases/surv.csv};
\addplot[omThm,very thick,dashed] table[x=year,y=survivors]{figdata/research/03-point-in-time-data-and-the-biases/surv.csv};
\legend{all listed names (with delistings),today's survivors only}
\end{semilogyaxis}
\end{tikzpicture}

\omcaption{Equal-weighted growth of one unit in a simulated market of 2\,000 stocks over twenty years, with 2.2\% of names
delisted each year at $-30\%$: the investable universe ends at 4.6, the universe of the names still listed at the end at 12.9.
Data: the chapter's module, seeded.}\label{fig:rs:point-in-time-data-and-the-biases:surv}
\end{omfigure}

\section{Restatements, revisions and vintages}

\begin{definition}[Restatement, data vintage, backfill bias]\label{def:rs:point-in-time-data-and-the-biases:vintage}
A \emph{restatement}\index{restatement} is a later published change to a value already released: a company restating an
earlier quarter's accounts, a statistical agency revising a figure. A \emph{data vintage}\index{data vintage} is the whole
history of a series as it was published at one \omterm{def:rs:point-in-time-data-and-the-biases:bitemporal}{knowledge time}. \emph{Backfill bias}\index{backfill bias} arises when history is
added to a database after the fact, typically when a fund or a vendor's new coverage enters with its past returns: the entrants
chose to enter because their past was good, and the backfilled history was never available in real time.
\end{definition}

The US payroll series shows how large revisions are for a headline number. The Philadelphia Fed's Real-Time Data Set for
Macroeconomists (Croushore and Stark, 2001) keeps every vintage of the series since 1964; vintage $m$ is the history as
published in the report of month $m$. For the 318 months from January 2000 to June 2026 with a first release, the monthly
change was later revised by 74\,000 jobs on average in absolute value, and by more than 100\,000 in 24\% of months; the
revisions average almost zero ($-3\,000$) but have a standard deviation of 106\,000. Eighteen months changed sign. The largest
revision is March 2020 ($-701\,000$ first, $-1\,398\,000$ today).

The September 2008 figure shows a single number's life (\cref{fig:rs:point-in-time-data-and-the-biases:life}): $-159\,000$,
then $-284\,000$ and $-403\,000$ in the next two reports, $-321\,000$ after the benchmark revision of February 2009, $-458\,000$
after that of February 2010, $-434\,000$ after that of February 2011, and $-451\,000$ today. The recession as a whole was invisible in its true
size while it happened (\cref{fig:rs:point-in-time-data-and-the-biases:y2008}): the first releases of the twelve months of 2008
sum to a loss of 1.88 million jobs, the January 2009 report put the year at 2.59 million, and today's data say 3.55 million.
A model fitted on today's data learns a relation between economic news and markets that no participant could have acted on.

\begin{omfigure}
\begin{tikzpicture}
\begin{axis}[width=10cm,height=4.8cm,xlabel={months after September 2008 (knowledge time)},ylabel={thousand jobs},
  tick label style={font=\small},label style={font=\small},xmin=0,xmax=37,ymin=-500,ymax=0,grid=major,grid style={gray!25},
  table/col sep=comma]
\addplot[omDef,very thick,const plot,mark=*,mark size=1pt] table[x=after,y=change]{figdata/research/03-point-in-time-data-and-the-biases/life.csv};
\addplot[black!60,thick,dashed,sharp plot] coordinates {(0,-451) (37,-451)};
\node[black!70,anchor=south,font=\footnotesize] at (axis cs:22.5,-446) {today: $-451$};
\end{axis}
\end{tikzpicture}

\omcaption{The US payroll change for September 2008 in each of the 36 monthly vintages after it: $-159$, $-284$, $-403$, then the
benchmark revisions of February 2009 ($-321$), February 2010 ($-458$) and February 2011 ($-434$); $-451$ in today's vintage
(dashed). Data: Federal Reserve Bank of Philadelphia, Real-Time Data Set for Macroeconomists (BLS data).}\label{fig:rs:point-in-time-data-and-the-biases:life}
\end{omfigure}

\begin{omfigure}
\begin{tikzpicture}
\begin{axis}[width=11cm,height=5cm,xlabel={month},ylabel={thousand jobs},tick label style={font=\small},
  label style={font=\small},xmin=-0.5,xmax=23.5,ymin=-1000,ymax=200,grid=major,grid style={gray!25},
  ytick={-1000,-500,0},yticklabels={$-1\,000$,$-500$,$0$},
  xtick={0,3,6,9,12,15,18,21},xticklabels={Jan 08,Apr 08,Jul 08,Oct 08,Jan 09,Apr 09,Jul 09,Oct 09},
  legend columns=2,legend style={font=\footnotesize,at={(0.5,-0.32)},anchor=north,draw=none,fill=none,
  /tikz/every even column/.append style={column sep=8pt}},table/col sep=comma]
\addplot[omThm,very thick,mark=*,mark size=1.2pt] table[x=i,y=first]{figdata/research/03-point-in-time-data-and-the-biases/y2008.csv};
\addplot[omDef,very thick,mark=square*,mark size=1.2pt] table[x=i,y=latest]{figdata/research/03-point-in-time-data-and-the-biases/y2008.csv};
\legend{first release,today's vintage}
\end{axis}
\end{tikzpicture}

\omcaption{Monthly US payroll changes in 2008--2009 as first published and as known today. The first releases understated
the losses of 2008 by 1.7 million jobs in total. Data: Federal Reserve Bank of Philadelphia, Real-Time Data Set for
Macroeconomists (BLS data).}\label{fig:rs:point-in-time-data-and-the-biases:y2008}
\end{omfigure}

Company data behave the same way. Accounts are restated; data vendors correct errors, change definitions and extend their
coverage backwards. A vendor that adds a thousand small companies with ten years of history has created a backfilled
universe, and a strategy backtested on it trades companies the firm could not have known it would cover. The cure is the same
in every case: keep the vintages.

\section{As-of joins and bitemporal storage}

\begin{definition}[As-of join]\label{def:rs:point-in-time-data-and-the-biases:asof}
An \emph{as-of join}\index{as-of join} matches each row of a left table (decision times) with the last row of a right table
(observations) whose time is at or before the left row's time, optionally within a tolerance; strictly before when the right
table's rows become known only after their stamp.
\end{definition}

An ordinary join on equal time stamps fails in both directions: it drops the decision times at which no observation has
exactly the same stamp, and it happily matches a decision with an observation stamped at the same instant that arrived later.
The \omterm{def:rs:point-in-time-data-and-the-biases:asof}{as-of join}, on \omterm{def:rs:point-in-time-data-and-the-biases:bitemporal}{knowledge times}, answers the only question a backtest may ask: what was the latest thing known?

\begin{method}[Making a backtest point-in-time]\label{met:rs:point-in-time-data-and-the-biases:pit}
\begin{enumerate}
\item Give every input a \omterm{def:rs:point-in-time-data-and-the-biases:bitemporal}{knowledge time}: the release or filing time, the arrival time of a message, the end of the \omterm{def:rs:market-data-for-research:bar}{bar} plus the
processing delay. Where only a \omterm{def:rs:point-in-time-data-and-the-biases:bitemporal}{valid time} exists, derive a conservative \omterm{def:rs:point-in-time-data-and-the-biases:bitemporal}{knowledge time} from a documented lag.
\item Store every vintage (\omterm{def:rs:point-in-time-data-and-the-biases:bitemporal}{bitemporal data}); never overwrite a value.
\item Read inputs through a guard that knows the decision time and refuses any later \omterm{def:rs:point-in-time-data-and-the-biases:bitemporal}{knowledge time}.
\item Build universes, identifiers and adjustments as of the decision time as well (chapter~4).
\item Test for leakage: perturb every value known after the decision time and check that no decision changes.
\end{enumerate}
\end{method}

\section{A catalogue of the other biases}

Look-ahead, survivorship and backfill are the biases of data; others come from the researcher. \emph{Selection of the sample
period}: starting a test after a crash, or ending it before one, is a trial the log must count (chapter~1). \emph{Mismatched
calendars}: a US close and a Tokyo close stamped with the same date are fifteen hours apart. \emph{Stale prices}: an illiquid
instrument's last trade may be days old, so a portfolio marked with it looks smoother than it is (chapter~22). \emph{Corporate
actions applied with hindsight}: a price series adjusted for a split with today's factor carries information about the split
into the past (chapter~4). \emph{Index membership}: testing a strategy on today's index members is \omterm{def:rs:point-in-time-data-and-the-biases:survivorship}{survivorship bias} under
another name. All are look-ahead in a broad sense, and all are prevented by the same discipline: at every decision time, only
what was known then.

\section{Tutorial: payrolls as they were known}\label{tut:rs:point-in-time-data-and-the-biases:payrolls}

\begin{tutorial}
\textbf{Goal.} Load the 2008--2009 payroll vintages into a \omterm{def:rs:point-in-time-data-and-the-biases:bitemporal}{bitemporal} store, ask what was known when, and let a guard refuse
the future. \textbf{End state:} \cref{fig:rs:point-in-time-data-and-the-biases:life,fig:rs:point-in-time-data-and-the-biases:y2008};
the 2008 total as known in January 2009 ($-2.59$ million) and today ($-3.55$ million).
\begin{tutsteps}
\item \textbf{The store.} Each (entity, field, \omterm{def:rs:point-in-time-data-and-the-biases:bitemporal}{valid time}) keeps a sorted list of (\omterm{def:rs:point-in-time-data-and-the-biases:bitemporal}{knowledge time}, value); a query finds the
last \omterm{def:rs:point-in-time-data-and-the-biases:bitemporal}{knowledge time} at or before the one asked.
\omcode{code/firm/pit/firm_pit.py}{39}{56}{Recording and reading a value by valid and knowledge time.}\label{lst:rs:point-in-time-data-and-the-biases:store}
\item \textbf{The guard.} A view of the store for one decision time; any query about later knowledge raises.
\omcode{code/firm/pit/firm_pit.py}{76}{90}{The look-ahead guard.}\label{lst:rs:point-in-time-data-and-the-biases:guard}
\item \textbf{Ask.} \texttt{real\_time\_view} for September 2008 with the decision time October 2008 returns $-159$, and with
September 2011 $-434$; a guard with the decision time November 2008 raises \texttt{LookAheadError} when asked for knowledge of
December 2008. Sum the twelve months of 2008 as known in each vintage and draw the result against the vintage date.
\end{tutsteps}
\textbf{What to change next.} Rebuild the data from the source with \texttt{rs\_fetch\_payrolls.py}; add the unemployment rate,
whose vintages the same data set provides, and find the months whose sign of change flipped.
\end{tutorial}

\section{Build: the point-in-time store}\label{bld:rs:point-in-time-data-and-the-biases:pit}

\begin{build}
\bfield{Purpose} Every fact the miniature firm uses in research is read through this store, so that a backtest sees only what
was known at each decision time. Chapter~4 builds the security master on it, chapter~11 the fundamental features.
\bfield{Interface} \texttt{Store.put(entity, field, valid, known, value)}; \texttt{asof(entity, field, valid, known)};
\texttt{first}, \texttt{vintages}, \texttt{snapshot(entity, field, known)}, \texttt{latest}; \texttt{Guard(store,
decision\_time)} with the same queries; \texttt{asof\_join(left\_t, right\_t, right\_v, tolerance, strict)}.
\bfield{Rules} Append-only by \omterm{def:rs:point-in-time-data-and-the-biases:bitemporal}{knowledge time} (a correction at an existing \omterm{def:rs:point-in-time-data-and-the-biases:bitemporal}{knowledge time} replaces it); times of any comparable
type used consistently; the guard raises rather than returning \texttt{None}, so a leak is loud.
\bfield{Acceptance tests} \texttt{code/firm/pit/tests/}: the life of the September 2008 figure; a snapshot equals a vintage and
\texttt{latest} equals today's view; the guard refuses a later \omterm{def:rs:point-in-time-data-and-the-biases:bitemporal}{knowledge time}; the \omterm{def:rs:point-in-time-data-and-the-biases:asof}{as-of join} with and without strictness and
tolerance.
\bfield{Stretch} A columnar version keyed by (entity, \omterm{def:rs:point-in-time-data-and-the-biases:bitemporal}{valid time}) for millions of rows; storing the \omterm{def:rs:point-in-time-data-and-the-biases:bitemporal}{knowledge time} of a
correction separately from that of the original value; an audit report of every query a backtest made.
\end{build}

\begin{omsources}
\item D.~Croushore and T.~Stark, ``A real-time data set for macroeconomists'', \emph{Journal of Econometrics} 105(1), 2001;
data: Federal Reserve Bank of Philadelphia, Real-Time Data Set for Macroeconomists, nonfarm payroll employment.
\item T.~Shumway, ``The delisting bias in CRSP data'', \emph{Journal of Finance} 52(1), 1997.
\item S.~J.~Brown, W.~Goetzmann, R.~G.~Ibbotson and S.~A.~Ross, ``Survivorship bias in performance studies'', \emph{Review of
Financial Studies} 5(4), 1992.
\item E.~J.~Elton, M.~J.~Gruber and C.~R.~Blake, ``Survivor bias and mutual fund performance'', \emph{Review of Financial Studies}
9(4), 1996.
\item W.~Fung and D.~A.~Hsieh, ``Performance characteristics of hedge funds and commodity funds: natural vs.\ spurious biases'',
\emph{Journal of Financial and Quantitative Analysis} 35(3), 2000.
\item R.~W.~Banz and W.~J.~Breen, ``Sample-dependent results using accounting and market data: some evidence'', \emph{Journal of
Finance} 41(4), 1986.
\end{omsources}

\section{Exercises}

\begin{exercise}[$\star$]\label{exo:rs:point-in-time-data-and-the-biases:1}
A daily \omterm{def:rs:market-data-for-research:bar}{bar} is stamped with its date and holds the close. A signal uses the close of day $t$ and the backtest trades at the close
of day $t$. What is wrong, and what is the earliest trade the signal allows?
\end{exercise}

\begin{exercise}[$\star$]\label{exo:rs:point-in-time-data-and-the-biases:2}
In a universe, 15\% of name-months belong to names that are later delisted, with an average return of $-20\%$ a year, and the
others average 10\%. By how much does a survivors-only average overstate the universe's?
\end{exercise}

\begin{exercise}[$\star$]\label{exo:rs:point-in-time-data-and-the-biases:3}
Using the vintages in \cref{fig:rs:point-in-time-data-and-the-biases:life}, what did a desk know about September 2008 in
January 2009, in June 2009 and in September 2011?
\end{exercise}

\begin{exercise}[$\star\star$]\label{exo:rs:point-in-time-data-and-the-biases:4}
Decision times are 10:00:00.0, 10:00:01.0 and 10:00:02.0; quotes are stamped 09:59:59.5, 10:00:01.0 and 10:00:01.8, and each
arrives 0.3 seconds after its stamp. Which quote does an \omterm{def:rs:point-in-time-data-and-the-biases:asof}{as-of join} on stamps give each decision, and which one on arrival
times?
\end{exercise}

\begin{exercise}[$\star\star$]\label{exo:rs:point-in-time-data-and-the-biases:5}
Payroll revisions from first release to today have a standard deviation of 106\,000 and a mean of $-3\,000$. A surprise model
reads the first release against a forecast whose error, in final data, has a standard deviation of 90\,000. What share of the
apparent variance of first-release surprises is revision noise, if the two are independent?
\end{exercise}

\begin{exercise}[$\star\star$]\label{exo:rs:point-in-time-data-and-the-biases:6}
A hedge-fund database adds funds with their past three years of returns when they join. Funds join when their past three years
beat the median. Explain the bias in the database's average return and how to remove it.
\end{exercise}

\begin{exercise}[$\star\star\star$]\label{exo:rs:point-in-time-data-and-the-biases:7}
\emph{Coding.} With \texttt{survivorship()}, measure the bias for barriers of $-1.2$ and $-2.0$ and for delisting returns of $0$
and $-100\%$. Which parameter moves the bias more, and why?
\end{exercise}

\begin{exercise}[$\star\star\star$]\label{exo:rs:point-in-time-data-and-the-biases:8}
\emph{Find the flaw.} ``We backtested our value strategy on the current members of the S\&P~500 since 2000, using each company's
fundamentals as they appear in our vendor's database today. It beats the index by 3\% a year.''
\end{exercise}

\section{Problem: The Survivors' Premium}

\begin{problem}[{Weekend problem --- how much does a list of today's stocks flatter the past?}]\label{pb:rs:point-in-time-data-and-the-biases:1}
The chapter's simulated market: 2\,000 names, random walks in log price with an 8\% drift and 40\% volatility a year, delisting
at a cumulative log return of $-1.6$ with a return of $-30\%$, replacement by a new listing, twenty years, seed 11.

\medskip\noindent\textbf{Part I --- The market.}
\begin{enumerate}
\item What price loss does a cumulative log return of $-1.6$ represent?
\item What fraction of names is delisted each year?
\item What is the annual equal-weighted return of the investable universe, delistings included?
\item And of today's survivors, over their own histories?
\item What is the bias, in points a year?
\end{enumerate}

\medskip\noindent\textbf{Part II --- Where it comes from.}
\begin{enumerate}[resume]
\item What share of name-months belongs to names that are later delisted?
\item What are $r_F$ and $r_S$?
\item What does \cref{prop:rs:point-in-time-data-and-the-biases:survivors} give, and why does it differ from Question 5?
\item What is the bias with a delisting return of zero?
\item What is it with a delisting return of $-100\%$?
\end{enumerate}

\medskip\noindent\textbf{Part III --- Sensitivity.}
\begin{enumerate}[resume]
\item What are the delisting rate and the bias with a barrier of $-1.2$?
\item And with a barrier of $-2.0$?
\item Why does a higher barrier (earlier delisting) raise the bias?
\item What would a survivors-only backtest of a strategy that buys stocks after large falls report?
\end{enumerate}

\medskip\noindent\textbf{Part IV --- Judgement.}
\begin{enumerate}[resume]
\item Where do you get the delisted names for a real backtest?
\item What delisting return do you use when a database has none, and why not zero?
\item How does the bias interact with an equal-weighted versus a capitalisation-weighted universe?
\item State the \emph{named result}: the annual return inflation of the survivors-only universe at the simulated delisting rate,
and the share of it that the delisting month explains.
\item Name two other biases with the same structure.
\item In one sentence: why is a list of today's companies a forecast?
\end{enumerate}
\end{problem}

\section{Interview questions}

\begin{interviewq}[$\star$ \iqroles{researcher, trader}]\label{iq:rs:point-in-time-data-and-the-biases:1}
What is \omterm{def:rs:point-in-time-data-and-the-biases:survivorship}{survivorship bias}, and how would it show up in a stock-selection backtest?
\end{interviewq}

\begin{interviewq}[$\star\star$ \iqroles{researcher}]\label{iq:rs:point-in-time-data-and-the-biases:2}
Your macro model uses GDP growth. Which GDP figure should it use for each date of the backtest?
\end{interviewq}

\begin{interviewq}[$\star\star$ \iqroles{developer, researcher}]\label{iq:rs:point-in-time-data-and-the-biases:3}
Design a table for fundamental data that supports point-in-time queries. What are its keys?
\end{interviewq}

\begin{interviewq}[$\star\star$ \iqroles{developer, mle}]\label{iq:rs:point-in-time-data-and-the-biases:4}
How would you test automatically that a feature pipeline has no look-ahead?
\end{interviewq}

\begin{interviewq}[$\star\star$ \iqroles{researcher, risk}]\label{iq:rs:point-in-time-data-and-the-biases:5}
A hedge-fund index shows high returns and low volatility. Which database biases would you suspect?
\end{interviewq}

\begin{interviewq}[$\star\star\star$ \iqroles{researcher}]\label{iq:rs:point-in-time-data-and-the-biases:6}
Show that a survivors-only average overstates a universe's average return by $w(r_S - r_F)$, and say which of the terms is larger
in practice.
\end{interviewq}
